\section{Conclusion}
\label{sec:conclusion}

% We introduced LexBench-Browser, a comprehensive benchmark addressing three systematic limitations in browser agent evaluation: task distribution shift, coarse-grained metrics, and fragmented infrastructure. Our dataset of 350+ tasks achieves balanced coverage of information-seeking and task-execution scenarios, our rubric-guided evaluation provides fine-grained diagnostic feedback, and our open infrastructure enables reproducible evaluation with authentication and anti-detection support.

% Experiments reveal that current agents exhibit substantial performance gaps on task-execution scenarios and authenticated workflows---capabilities essential for practical utility but underrepresented in existing benchmarks. Our evaluation methodology achieves higher agreement with human judgments while reducing computational costs.

% \textbf{Limitations.} Our benchmark focuses on Chinese and English websites, leaving other languages for future work. The dynamic nature of web environments means some tasks may require periodic updates. Our safety evaluation, while more realistic than static approaches, cannot cover all possible threat vectors.


% LexBench-Browser highlights a practical gap that is easy to miss in today's web agent evaluations. In our experiments, success on T2 task execution remains consistently low, with a persistent gap from T1 information retrieval. This suggests that current agents are much better at observing and extracting information than at reliably taking state-changing actions such as authenticated flows, form filling, and multi-step transactions. A key factor is context rot: as T2 tasks require longer interaction sequences, the agent's context accumulates stale observations, irrelevant DOM fragments, and compounding errors, degrading decision quality over time.
% At the same time, agents do not show a clear advantage on T1 tasks. For many retrieval-heavy slices, text-based Deep Research style systems can be competitive—sometimes even stronger—because they operate on cleaner textual context without the noise and fragility of GUI interaction. This implies that a browser-wrapped agent is not automatically the best tool when the job is primarily information seeking.
% Safety is another bottleneck. By evaluating on live websites with dynamic safety scenarios, we observe that current agents can still be steered into unsafe actions or continue along harmful workflows, which prevents users from confidently delegating control.
% Finally, latency and token consumption are major constraints. Architectural overhead in some frameworks significantly increases per-task cost. LexBench-Browser surfaces these issues under realistic task distributions and provides diagnostic signals that make improvement actionable: close the T2 execution gap, reach a trustworthy safety bar, and reduce end-to-end cost.
LexBench-Browser reveals a clear and persistent gap between information retrieval and task execution. While current agents achieve reasonable performance on T1 tasks, their success on T2 remains consistently low, showing that they are much better at observing and extracting information than at reliably executing state-changing workflows such as authenticated interactions, form filling, and multi-step transactions. This weakness is closely related to long-horizon interaction issues, where accumulated context noise and compounding errors gradually degrade decision quality. Moreover, agents do not show a decisive advantage even on T1 tasks, as text-based Deep Research style systems can often be competitive due to their cleaner and more stable textual context. Together with the remaining safety vulnerabilities and the high latency and token cost of current frameworks, these results indicate that today’s browser agents are not yet a dependable drop-in solution. LexBench-Browser makes these limitations explicit under realistic task distributions and provides actionable diagnostic signals, pointing to three concrete directions for progress: closing the T2 execution gap, reaching a trustworthy safety bar, and reducing end-to-end cost.